\chapter{Diseño y Justificación}
\label{ch:diseno-justificacion}

Este capítulo presenta el diseño detallado del \textbf{CMOS Microscopy Payload
(CMOSM)} y justifica las decisiones adoptadas en los dominios óptico, mecánico,
electrónico y de software. Las elecciones de diseño se fundamentan en los
requerimientos definidos en el Capítulo~\ref{ch:requerimientos}, los modos de
operación descritos en el Capítulo~\ref{ch:modos-operacion} y la arquitectura
funcional presentada en el Capítulo~\ref{ch:funcional}.

El enfoque de diseño prioriza la simplicidad, reproducibilidad experimental y
trazabilidad de datos, alineándose con el carácter académico del proyecto
INTISAT y las restricciones propias de un CubeSat.

% -------------------------------------------------------------------
\section{Justificación del enfoque de microscopía}

El sistema CMOSM adopta un enfoque de \textbf{microscopía sin lentes en
configuración de contacto}, en lugar de una arquitectura óptica convencional
basada en lentes. Esta decisión se justifica por:

\begin{itemize}
    \item Reducción significativa de complejidad mecánica y óptica.
    \item Eliminación de elementos refractivos sensibles a desalineación.
    \item Compatibilidad con arquitecturas compactas de bajo volumen y masa.
    \item Adecuación para demostradores tecnológicos y misiones educativas.
\end{itemize}

Los ensayos experimentales realizados demuestran que, bajo condiciones adecuadas
de iluminación y acoplamiento óptico, la microscopía de contacto permite la
identificación de estructuras celulares relevantes en muestras biológicas
planas, cumpliendo el objetivo científico del payload.

\begin{figure}[H]
    \centering
    \includegraphics[width=0.7\textwidth]{figures/Captura_muestra1.png}
    \caption{Captura de tejido vegetal mediante microscopía de contacto sin lentes, mostrando estructuras celulares identificables.}
    \label{fig:captura_muestra1}
\end{figure}

\begin{figure}[H]
    \centering
    \includegraphics[width=0.7\textwidth]{figures/Captura_muestra2.jpg}
    \caption{Captura de muestra biológica con iluminación optimizada, demostrando el contraste alcanzable con el sistema CMOSM.}
    \label{fig:captura_muestra2}
\end{figure}

% -------------------------------------------------------------------
\section{Diseño óptico y justificación}

\subsection{Sensor CMOS}

El sensor CMOS OV5647 fue seleccionado por su disponibilidad como componente
\emph{Commercial Off-The-Shelf} (COTS), su resolución nativa de 2592×1944 píxeles
y su compatibilidad con plataformas embebidas como la Raspberry Pi.

La operación del sensor se realiza bajo control manual de exposición y ganancia,
deshabilitando mecanismos automáticos de ajuste. Esta decisión permite:

\begin{itemize}
    \item Consistencia entre capturas consecutivas.
    \item Control explícito de la integración fotónica.
    \item Reproducibilidad experimental.
\end{itemize}

\subsection{Iluminación}

El sistema de iluminación emplea una pantalla OLED de 1.3 pulgadas controlada vía
bus I2C como fuente plana y programable. Esta solución se eligió frente a arreglos
de LEDs discretos por:

\begin{itemize}
    \item Flexibilidad para generar múltiples patrones sin modificar hardware.
    \item Integración compacta con bajo consumo de potencia.
    \item Capacidad de conmutación rápida de patrones.
\end{itemize}

En operación nominal se utilizan patrones uniformes para maximizar el contraste
de la muestra. La iluminación estructurada se reserva para modos experimentales.

\subsection{Interfaz muestra--sensor}

La interfaz entre la muestra y el sensor se reconoce como un elemento crítico
del desempeño óptico del sistema. Durante la fase actual del proyecto, se
emplean soluciones experimentales de laboratorio para validar el principio de
acoplamiento óptico y su impacto en la calidad de imagen.

La selección e implementación de un material definitivo para esta interfaz se
identifica como trabajo futuro y no forma parte del alcance del presente CDR.

% -------------------------------------------------------------------
\section{Diseño mecánico y justificación}

El diseño mecánico del CMOSM prioriza la simplicidad estructural y la estabilidad
geométrica del sistema. La arquitectura mecánica se concibe como un soporte
rígido que mantiene:

\begin{itemize}
    \item Posicionamiento estable del sensor CMOS.
    \item Alineamiento fijo entre el sensor y el sistema de iluminación.
    \item Acceso controlado a la superficie del sensor para preparación de
    muestras.
\end{itemize}

Esta aproximación reduce la necesidad de mecanismos de ajuste fino y minimiza
riesgos asociados a vibraciones y desalineaciones durante operación de
laboratorio.

% -------------------------------------------------------------------
\section{Diseño electrónico y justificación}

Desde el punto de vista electrónico, el CMOSM se apoya en la Raspberry Pi como
computadora de carga útil. Esta elección se justifica por:

\begin{itemize}
    \item Compatibilidad directa con el sensor OV5647.
    \item Disponibilidad de librerías maduras (Picamera2).
    \item Capacidad de procesamiento suficiente para control y análisis básico
    de imágenes.
\end{itemize}

El control del OLED mediante bus I2C permite una implementación simple y robusta,
con bajo número de líneas y reducido consumo de potencia, alineado con los
requerimientos eléctricos del payload.

Se incluyen consideraciones básicas de protección eléctrica, tales como manejo
adecuado de niveles lógicos y prevención de descargas electrostáticas en
conexiones externas.

% -------------------------------------------------------------------
\section{Diseño de software y justificación}

El software del CMOSM se diseña bajo una arquitectura modular, separando
claramente las funciones de control, adquisición, procesamiento y almacenamiento
de datos.

Las decisiones clave de diseño de software incluyen:

\begin{itemize}
    \item Implementación de una máquina de estados para gestión de modos
    operativos.
    \item Control manual explícito de parámetros críticos del sensor.
    \item Registro automático de metadatos para cada captura.
    \item Inclusión de rutinas de calibración básica (dark frame y flat field).
\end{itemize}

El procesamiento computacional avanzado, incluyendo reconstrucción FPM y
super--resolución basada en aprendizaje profundo, se implementa como extensión
experimental, sin afectar la operación nominal del sistema.

% -------------------------------------------------------------------
\section{Justificación del uso de técnicas computacionales}

Las técnicas de reconstrucción computacional y realce visual se incorporan con
un enfoque conservador y evaluativo. En particular:

\begin{itemize}
    \item Fourier Ptychographic Microscopy (FPM) se emplea como herramienta de
    investigación experimental.
    \item La super--resolución basada en redes neuronales se utiliza únicamente
    como realce visual, sin introducir información física nueva.
\end{itemize}

\begin{figure}[H]
    \centering
    \includegraphics[width=0.8\textwidth]{figures/fpm_reconstruction1.png}
    \caption{Ejemplo de reconstrucción mediante Fourier Ptychographic Microscopy (FPM), mostrando el potencial de las técnicas computacionales para la mejora de resolución.}
    \label{fig:fpm_reconstruction}
\end{figure}

Esta distinción permite aprovechar las ventajas de dichas técnicas sin
comprometer la validez física ni la trazabilidad de los resultados obtenidos por
el sistema.

% -------------------------------------------------------------------
\section{Resumen del diseño}

El diseño del CMOSM refleja un compromiso consciente entre simplicidad,
funcionalidad y capacidad experimental. Las decisiones adoptadas permiten validar
el sistema como demostrador tecnológico de microscopía sin lentes, manteniendo
una arquitectura extensible para futuras iteraciones del proyecto INTISAT.
